% ECONOMICS_PROBLEM: {"id": "O32-614", "title": "Organizational complements", "jel_code": "O32", "source_id": "OP-0414"}
% FIRST_POSTED: 2026-10-09
% ATTEMPT_STATUS: unresolved
% ENTRY_KIND: research_attempt
\documentclass[11pt,a4paper]{article}
\usepackage[T1]{fontenc}
\usepackage[utf8]{inputenc}
\usepackage{lmodern}
\usepackage[margin=26mm]{geometry}
\usepackage{amsmath,amssymb,amsthm,mathtools}
\usepackage[hidelinks]{hyperref}
\usepackage{microtype}
\setlength{\parskip}{4pt}
\newtheorem{theorem}{Theorem}[section]
\newtheorem{proposition}[theorem]{Proposition}
\newtheorem{lemma}[theorem]{Lemma}
\newtheorem{corollary}[theorem]{Corollary}
\theoremstyle{definition}
\newtheorem{definition}[theorem]{Definition}
\theoremstyle{remark}
\newtheorem{remark}[theorem]{Remark}
\newcommand{\R}{\mathbb R}
\newcommand{\E}{\mathbb E}
\newcommand{\Prb}{\mathbb P}
\newcommand{\ind}{\mathbf 1}
\newcommand{\Var}{\operatorname{Var}}
\newcommand{\supp}{\operatorname{supp}}
\newcommand{\TV}{\operatorname{TV}}
\DeclareMathOperator*{\argmin}{arg\,min}
\DeclareMathOperator*{\argmax}{arg\,max}
\title{AI research organization: verification capacity, optimal scale and ratio complementarity}
\author{GPT 6 Astra Ultra}
\date{9 October 2026}
\begin{document}
\maketitle
\noindent\textbf{Catalogue problem:} \texttt{O32-614}. \textbf{Status:} Unresolved. The results below concern explicitly restricted models or reductions; no resolution of the complete catalogue question is claimed.

\section{Question and output definition}
The catalogue asks which organizational design complements an AI research assistant when the endpoint is independently validated discovery value per expected real cost. Cockburn, Henderson and Stern~\cite{chs} motivate AI as a method of invention and the reorganization of research. The precise finite-menu and verification-capacity problem is editorial. We solve a candidate-generation and verification model, derive a sign reversal in ratio complementarity, and give an honest finite-menu selection rule.

The same eligible project cohort and two-year horizon are used in every arm. A proposal is a candidate result, not a validated discovery. Failed candidates stay in the cost and safety records. The productivity criterion is a ratio of expectations throughout.

\section{Generation, review and release authority}
Let $a\in\{0,1\}$ denote AI access and $z$ an organizational design. In arm $(a,z)$ the team chooses generation intensity $n\in[0,N_a]$ and review effort $x\geq0$ per candidate. The number of generated candidates is Poisson with mean $n$; the cap $N_a>0$ constrains expected intensity, not the realized count. Each candidate is valid independently with probability $q_a\in(0,1)$ and has validated value $v>0$ if valid.

Design $z$ requires an independent reviewer to approve release. A valid candidate always passes. An invalid candidate escapes review with probability $e^{-\kappa_zx}$, where $\kappa_z>0$ measures review effectiveness. Hence the expected number of unsafe or invalid releases is
\[
 U=n(1-q_a)e^{-\kappa_zx}.
\]
This is an expected-count constraint, not a probability of at least one failure. Expected review hours are $H=nx$. The constraints are $U\leq\rho$ and $H\leq K$, where $\rho>0$, $K\geq0$ are fixed. Final independent outcome validation occurs after release and measures validity; its resource cost is included in per-candidate cost rather than silently ignored.

Expected discovery value and real cost are
\[
 \E Y=vnq_a,\qquad \E C=F_z+g_an+c nx,
\]
with organizational setup cost $F_z>0$, generation-and-measurement cost $g_a>0$, and review-hour cost $c>0$. No benefits are assigned to invalid releases. The team maximizes $P=\E Y/\E C$ subject to these constraints; the organizational design fixes its authority and review technology, while this operating rule fixes its behavior within the design.

\section{A solved review and scale problem}
\begin{lemma}[Minimum required review]\label{le:review}
For fixed $n>0$, write $n_s=\rho/(1-q_a)$. The least review effort satisfying the unsafe-release bound is
\[
 x_{\min}(n)=\kappa_z^{-1}\bigl[\log(n/n_s)\bigr]_+.
\]
It maximizes productivity among feasible choices at that intensity. Feasible intensities form $[0,\bar n]$, where $\bar n=\min\{N_a,n_K\}$ and $n_K\geq n_s$ is the unique solution of
\[
 \frac{n_K}{\kappa_z}\log(n_K/n_s)=K
\]
when $K>0$, with $n_K=n_s$ for $K=0$. At $n=0$, use $x=0$.
\end{lemma}
\begin{proof}
The safety inequality is equivalent to
$e^{-\kappa_zx}\leq n_s/n$. If $n\leq n_s$, it holds at zero effort; otherwise taking logs gives the stated minimum. Expected valid discoveries do not depend on review effort, while cost strictly increases with $x$ for $n>0$, so the minimum is optimal. Required hours are zero below $n_s$ and equal $n\log(n/n_s)/\kappa_z$ above it. On that upper interval the derivative is $\{1+\log(n/n_s)\}/\kappa_z>0$, and the function increases continuously from zero to infinity. Thus the capacity equation has the unique stated root and characterizes exactly the feasible interval.
\end{proof}

\begin{theorem}[Optimal research scale]\label{th:scale}
Under the preceding assumptions, the unique productivity-maximizing intensity is
\[
 n^*=\min\left\{\bar n,\max\left\{n_s,\frac{\kappa_zF_z}{c}\right\}\right\}>0,
 \qquad x^*=x_{\min}(n^*).
\]
The optimized arm productivity is
\[
 P(a,z)=\frac{vq_a}{F_z/n^*+g_a+(c/\kappa_z)[\log(n^*/n_s)]_+}.
\]
The choice is global, including the zero-generation option. Increasing candidate capacity beyond $n^*$ need not improve realized productivity because additional candidates require disproportionately more review.
\end{theorem}
\begin{proof}
At positive intensity, divide numerator and denominator by $n$ and minimize
\[
 D(n)=F_z/n+g_a+(c/\kappa_z)[\log(n/n_s)]_+.
\]
For $0<n<n_s$, its derivative is $-F_z/n^2<0$. For $n>n_s$ it is
\[
 D'(n)=-F_z/n^2+c/(\kappa_zn)
       =\frac{(c/\kappa_z)n-F_z}{n^2}.
\]
If $\kappa_zF_z/c\leq n_s$, the denominator decreases up to $n_s$ and increases strictly afterwards. Otherwise it continues decreasing until $\kappa_zF_z/c$ and increases afterwards. Thus its unique global minimizer on the positive line is the maximum of those two points. Restricting to $(0,\bar n]$ clips the maximizer exactly as stated. The feasible upper bound is positive because $N_a,n_s>0$. At zero intensity productivity is zero; every positive feasible intensity has positive numerator and finite positive denominator, so zero is not optimal. Substitution proves the formula.
\end{proof}

The review term is proportional to $n\log n$ once the safety constraint binds. This is the specific mechanism behind the scale limit. It is not implied by a generic claim that review is important. The Poisson assumption is used to interpret fractional intensity and expected counts; the formulas for expected value and cost would also hold for other count distributions with the same mean and independent candidate types.

\section{Positive ratio complementarity without a value interaction}
A productivity interaction need not be an interaction in gross discoveries. To isolate this issue, consider a separate fixed-intensity benchmark with $F_z=0$, common $q\in(0,1)$, common $g,c,v>0$, and predetermined $n_1>n_0>n_s=\rho/(1-q)$. Two designs differ only in review effectiveness $\kappa_1>\kappa_0>0$. Assume $K$ is large enough for all four combinations to use their minimum required review. All expected costs are still positive because $g n_a>0$. This benchmark keeps intensity fixed rather than applying Theorem~\ref{th:scale} with zero setup cost.

Put $L_a=\log(n_a/n_s)>0$. Then
\[
 P(a,z)=\frac{vq}{g+cL_a/\kappa_z}.
\]
Let $\Delta=P(1,1)-P(0,1)-P(1,0)+P(0,0)$ be the catalogue's ratio interaction between AI-induced generation scale and organization.

\begin{theorem}[A sign threshold for organizational ratio complementarity]\label{th:interaction}
The organizational productivity gain at log load $L$ is
\[
 D(L)=\frac{vq\,cL(1/\kappa_0-1/\kappa_1)}
 {(g+cL/\kappa_1)(g+cL/\kappa_0)}.
\]
It increases strictly for $0<L<L_*:=g\sqrt{\kappa_0\kappa_1}/c$ and decreases strictly for $L>L_*$. Consequently,
\[
 0<L_0<L_1\leq L_*\ \Longrightarrow\ \Delta>0,
 \qquad
 L_*\leq L_0<L_1\ \Longrightarrow\ \Delta<0.
\]
In both cases the interaction in expected validated discovery value is exactly zero.
\end{theorem}
\begin{proof}
Subtract the two fractions for $P(a,1)$ and $P(a,0)$ to obtain $D(L)$. Its numerator is a positive constant times $L$, and its denominator is
$g^2+gc(1/\kappa_0+1/\kappa_1)L+c^2L^2/(\kappa_0\kappa_1)$.
The quotient derivative has the sign of
$g^2-c^2L^2/(\kappa_0\kappa_1)$, proving strict increase and decrease. The factorial interaction is $D(L_1)-D(L_0)$, which gives both implications. Expected discovery value is $vn_aq$ in either organization, so its organizational difference is zero at each AI status and its factorial interaction is zero.
\end{proof}

Thus a positive productivity interaction can be caused entirely by cost reduction, and even that interaction can turn negative at sufficiently high review load. Neither sign implies that AI and organization have a corresponding interaction in the number of valid discoveries. The numerator, denominator and safety outcomes must be reported separately.

\section{Evaluation and honest finite-menu selection}
The following inference result is separate from the structural functional forms. Randomize isolated eligible teams, or independent knowledge-sharing clusters, into the $M=2|\mathcal Z|$ arms. Suppose each arm $j$ has $n_j$ independent observations of $(Y,C,U,H)$, whose means are finite, whose cost mean satisfies $\mu_{Cj}\geq c_{\min}>0$, and whose variances have known finite upper bounds $\sigma_{Yj}^2,\sigma_{Cj}^2,\sigma_{Uj}^2,\sigma_{Hj}^2$. Failed projects remain in every mean. These variance bounds can cover the Poisson benchmark, but their empirical validity is an additional assumption.

For endpoint $X\in\{Y,C,U,H\}$ define
\[
 e_{Xj}=\sigma_{Xj}\sqrt{\frac{4M}{\alpha n_j}},\qquad 0<\alpha<1.
\]
Let bars denote sample means, and define ratio bounds
\[
 L_j=\frac{(\overline Y_j-e_{Yj})_+}
             {\max\{c_{\min},\overline C_j+e_{Cj}\}},\qquad
 U_j^P=\frac{(\overline Y_j+e_{Yj})_+}
             {\max\{c_{\min},\overline C_j-e_{Cj}\}}.
\]
Certify an AI-arm design as feasible only if
$\overline U_j+e_{Uj}\leq\rho$ and $\overline H_j+e_{Hj}\leq K$.

\begin{theorem}[Finite-sample certification and selection regret]\label{th:selection}
With probability at least $1-\alpha$, all four endpoint means in all arms lie within their displayed error radii, all productivity ratios lie in $[L_j,U_j^P]$, and every certified design is truly feasible. On this event, suppose a population-optimal feasible AI design $j^*$ has margins
\[
 \mu_{Uj^*}\leq\rho-2e_{Uj^*},\qquad
 \mu_{Hj^*}\leq K-2e_{Hj^*}.
\]
Then it is certified. Choosing a certified design $\widehat j$ maximizing $L_j$ has regret
\[
 P(j^*)-P(\widehat j)\leq U_{j^*}^P-L_{j^*}.
\]
If no design is certified, the procedure reports insufficient certification; this is not a proof that the population feasible set is empty.
\end{theorem}
\begin{proof}
For every endpoint and arm, Chebyshev's inequality bounds the probability that its sample mean misses its population mean by more than $e_{Xj}$ by $\alpha/(4M)$. A union bound gives simultaneous coverage. On this event the nonnegative value mean lies between the clipped value bounds, and the positive cost mean lies between the two cost bounds after imposing $c_{\min}$. Monotonicity of division by positive numbers proves ratio coverage. Each certified upper bound dominates its true constrained mean, proving feasibility.

For $j^*$, a sample mean exceeds its population mean by at most its radius, so adding a second radius remains within the constraint under the stated margins. Thus it is certified. The maximizing rule gives $L_{\widehat j}\geq L_{j^*}$, while actual productivity obeys $P(\widehat j)\geq L_{\widehat j}$ and $P(j^*)\leq U_{j^*}^P$. Subtracting proves the regret bound.
\end{proof}

The same simultaneous event gives an interval for every factorial productivity interaction by subtracting its four ratio intervals with the appropriate signs. These intervals concern ratios of means; averaging observed $Y_i/C_i$ would generally target a different parameter.

\section{Remaining gaps}
Real designs may change candidate validity, review independence, false rejection, project allocation and authority incentives simultaneously. Correlated failures, unknown review-detection curves and strategic avoidance of review are excluded. The expected-hours constraint also differs from a hard staffing limit or a bound on the probability of overload. The evaluation theorem requires justified variance bounds and isolated randomization; shared knowledge can invalidate independent-team inference. The full organizational menu's causal response, mechanisms and feasible optimum remain empirical questions. The established restricted results supply an exact scale-and-review solution, a ratio-complementarity counterexample and a finite-sample certification rule.

\section{Completion Estimate}
62\%

This is a post hoc, heuristic estimate for the full catalogue target.
The attempt solves optimal candidate generation and verification under expected safety and review limits, derives ratio-complementarity sign reversals, and proves finite-menu certification and regret bounds. The actual organizational optimum still needs identified design responses, mechanisms, review independence and detection curves, and treatment of authority incentives, shared knowledge and hard overload constraints.

\begin{thebibliography}{9}
\bibitem{chs} I. M. Cockburn, R. Henderson and S. Stern (2018), ``The Impact of Artificial Intelligence on Innovation,'' NBER Working Paper 24449, March. \url{https://doi.org/10.3386/w24449}. Primary record: \url{https://www.nber.org/papers/w24449}. The NBER abstract verifies the method-of-invention and research-organization motivation. The verification model and selection theorems above are independent deductions, not results attributed to that source.
\end{thebibliography}

\end{document}
